\documentclass[11pt, a4paper]{article}
% ══════════════════════════════════════════════════════════════
% Preâmbulo compartilhado pelos documentos da P2
% Reaproveita o estilo de docs/documento-unificado.tex
% ══════════════════════════════════════════════════════════════

\usepackage[utf8]{inputenc}
\usepackage[T1]{fontenc}
\usepackage[brazil]{babel}

\usepackage[top=2.5cm, bottom=2.5cm, left=3cm, right=2cm]{geometry}
\usepackage{setspace}
\onehalfspacing
\setlength{\parindent}{0pt}
\setlength{\parskip}{6pt}

\usepackage{lmodern}
\usepackage{microtype}
\usepackage{amssymb}
\usepackage{amsmath}
\usepackage{pifont}
\newcommand{\cmark}{\ding{51}}
\newcommand{\xmark}{\ding{55}}

\usepackage[table,dvipsnames]{xcolor}
\definecolor{accent}{HTML}{2C7A4B}
\definecolor{accentdim}{HTML}{1E5934}
\definecolor{accentlight}{HTML}{E8F4ED}
\definecolor{tabheader}{HTML}{2C3E50}
\definecolor{tabrow}{HTML}{F5F6FA}
\definecolor{linkblue}{HTML}{2980B9}
\definecolor{codebg}{HTML}{F5F5F5}
\definecolor{codeborder}{HTML}{D0D0D0}
\definecolor{ndvilow}{HTML}{8B4513}
\definecolor{ndvimid}{HTML}{DAA520}
\definecolor{ndvihigh}{HTML}{006400}
\definecolor{warnyellow}{HTML}{F39C12}
\definecolor{successgreen}{HTML}{27AE60}
\definecolor{errorred}{HTML}{E74C3C}

\usepackage{booktabs}
\usepackage{tabularx}
\usepackage{multirow}
\usepackage{makecell}
\usepackage{float}
\usepackage{graphicx}
\usepackage{caption}
\usepackage{array}

\usepackage{enumitem}
\setlist[itemize]{itemsep=2pt, parsep=0pt, topsep=4pt, leftmargin=*}
\setlist[enumerate]{itemsep=2pt, parsep=0pt, topsep=4pt, leftmargin=*}

\usepackage[colorlinks=true, linkcolor=linkblue, urlcolor=linkblue, citecolor=linkblue]{hyperref}

\usepackage{fancyhdr}
\pagestyle{fancy}
\fancyhf{}
\fancyhead[L]{\small\textcolor{gray}{EasyGis}}
\fancyhead[R]{\small\textcolor{gray}{P2 — Gerência e Manutenção de Software · Prof.\ Mauro · 2026.1}}
\fancyfoot[C]{\thepage}
\renewcommand{\headrulewidth}{0.4pt}

\usepackage{titlesec}
\titleformat{\section}
  {\Large\bfseries\color{accentdim}}
  {\thesection.}{0.5em}{}
\titleformat{\subsection}
  {\large\bfseries\color{accent}}
  {\thesubsection}{0.5em}{}
\titleformat{\subsubsection}
  {\normalsize\bfseries}
  {\thesubsubsection}{0.5em}{}

% Listings (código)
\usepackage{listings}
\lstdefinestyle{shellstyle}{
  basicstyle=\ttfamily\small,
  backgroundcolor=\color{codebg},
  frame=single,
  rulecolor=\color{codeborder},
  framesep=4pt,
  xleftmargin=8pt,
  xrightmargin=4pt,
  breaklines=true,
  showstringspaces=false,
  keepspaces=true,
  literate={─}{{-}}1 {├}{{|}}1 {│}{{|}}1 {└}{{`}}1
}
\lstdefinestyle{pythonstyle}{
  basicstyle=\ttfamily\footnotesize,
  language=Python,
  backgroundcolor=\color{codebg},
  frame=single,
  rulecolor=\color{codeborder},
  framesep=4pt,
  xleftmargin=8pt,
  xrightmargin=4pt,
  breaklines=true,
  keywordstyle=\color{accent}\bfseries,
  stringstyle=\color{ndvimid},
  commentstyle=\color{gray}\itshape,
  showstringspaces=false,
  keepspaces=true
}
\lstset{style=shellstyle}

% Caixas para código inline
\newcommand{\code}[1]{\texttt{\colorbox{codebg}{\strut#1}}}

% Caixas de destaque
\usepackage{tcolorbox}
\tcbuselibrary{skins, breakable}

\newtcolorbox{infobox}[1][]{%
  colback=accentlight, colframe=accent,
  boxrule=0.6pt, arc=2pt,
  left=8pt, right=8pt, top=4pt, bottom=4pt,
  fonttitle=\bfseries, breakable, #1
}

\newtcolorbox{warnbox}[1][]{%
  colback=yellow!8, colframe=warnyellow,
  boxrule=0.6pt, arc=2pt,
  left=8pt, right=8pt, top=4pt, bottom=4pt,
  breakable, #1
}

\newtcolorbox{successbox}[1][]{%
  colback=green!4, colframe=successgreen,
  boxrule=0.6pt, arc=2pt,
  left=8pt, right=8pt, top=4pt, bottom=4pt,
  breakable, #1
}


\title{%
  \vspace{-1.5cm}
  {\Huge\bfseries\color{accentdim} Plano de Testes}\\[8pt]
  {\Large\bfseries EasyGis}\\[18pt]
  {\normalsize\itshape Avaliação P2 --- Engenharia de Software, 2026.1}
}
\author{Equipe EasyGis}
\date{Versão 1.0 --- 09 de maio de 2026}

\begin{document}
\maketitle
\thispagestyle{empty}

\vfill
\begin{center}
\small\textit{Instituto de Ensino Superior --- ICEV}\\
\small\textit{Teresina, PI, Brasil}
\end{center}
\newpage

\tableofcontents
\newpage

% ══════════════════════════════════════════════════════════════
\section{Objetivo}
% ══════════════════════════════════════════════════════════════

Definir a estratégia, escopo e casos de teste para validar o MVP do sistema, com foco nos requisitos funcionais críticos de processamento de imagens Sentinel-2 (NDVI, EVI, estatísticas) e nos endpoints HTTP da API.

% ══════════════════════════════════════════════════════════════
\section{Escopo}
% ══════════════════════════════════════════════════════════════

\subsection{No escopo}

\begin{itemize}
  \item \textbf{Cálculo de índices espectrais} (NDVI, EVI) --- funções puras do \texttt{SentinelProcessor}
  \item \textbf{Estatísticas e histograma} --- agregações sobre arrays NumPy
  \item \textbf{Endpoints HTTP} --- \texttt{/}, \texttt{/health}, \texttt{/products}, validação de entrada de \texttt{/calculate-index}
  \item \textbf{Conversão de coordenadas} --- geração de polígonos Shapely a partir de coordenadas geográficas
  \item \textbf{Tratamento de casos de borda} --- divisão por zero, NaN, arrays vazios
\end{itemize}

\subsection{Fora do escopo}

\begin{itemize}
  \item \textbf{Integração com produto Sentinel-2 real} --- exigiria download de aproximadamente 1\,GB de dados; testado manualmente na demonstração.
  \item \textbf{Testes de UI} (frontend Next.js) --- fora do escopo da entrega P2 acordado com o professor; testes manuais via demonstração.
  \item \textbf{Performance e carga} --- sistema é local, single-user; não há requisito de carga.
  \item \textbf{Endpoints \texttt{/api/experiments/run} e \texttt{/api/classification/classify}} --- exigem produto Sentinel-2 real para teste end-to-end. Validação por inspeção manual.
\end{itemize}

% ══════════════════════════════════════════════════════════════
\section{Estratégia de Teste}
% ══════════════════════════════════════════════════════════════

\subsection{Níveis}

\begin{table}[H]
\centering
\small
\rowcolors{2}{tabrow}{white}
\begin{tabularx}{\textwidth}{@{}llX@{}}
\toprule
\rowcolor{tabheader}\textcolor{white}{\textbf{Nível}} & \textcolor{white}{\textbf{Ferramenta}} & \textcolor{white}{\textbf{Cobertura Pretendida}}\\
\midrule
Unitário & pytest & Funções puras de \texttt{sentinel\_processor.py} (cálculo de índices, estatísticas)\\
Integração & pytest + \texttt{TestClient} & Endpoints HTTP, validação Pydantic\\
Manual & Demonstração ao vivo & Fluxo completo end-to-end com produto Sentinel-2 e KML real\\
\bottomrule
\end{tabularx}
\end{table}

\subsection{Abordagem}

\begin{itemize}
  \item \textbf{Caixa branca} nas funções de cálculo: validação contra resultados matemáticos conhecidos.
  \item \textbf{Caixa preta} nos endpoints: requisições HTTP com payloads válidos/inválidos, verificação de status code e estrutura de resposta.
  \item \textbf{Fixtures sintéticas} (NumPy) substituem o I/O de imagens reais, permitindo execução rápida ($<1$\,s) e reprodutível.
\end{itemize}

\subsection{Critérios de aceitação}

\begin{itemize}
  \item 100\% dos casos de teste planejados executados.
  \item Taxa de sucesso $\geq 95\%$.
  \item Falhas analisadas e documentadas no Relatório de Testes (\texttt{09-relatorio-de-testes.pdf}).
\end{itemize}

% ══════════════════════════════════════════════════════════════
\section{Ambiente de Teste}
% ══════════════════════════════════════════════════════════════

\begin{table}[H]
\centering
\small
\rowcolors{2}{tabrow}{white}
\begin{tabularx}{\textwidth}{@{}lX@{}}
\toprule
\rowcolor{tabheader}\textcolor{white}{\textbf{Item}} & \textcolor{white}{\textbf{Especificação}}\\
\midrule
Sistema operacional & macOS 15 (Darwin 25.4.0)\\
Python & 3.12.13\\
pytest & 9.0.3\\
Gerenciador de pacotes & uv 0.11.3\\
Dependências & \texttt{pyproject.toml} grupo \texttt{[dependency-groups.dev]}\\
\bottomrule
\end{tabularx}
\end{table}

\textbf{Comando de execução:}

\begin{lstlisting}[style=shellstyle]
uv sync --group dev
uv run pytest tests/ -v --cov=api --cov-report=term-missing
\end{lstlisting}

% ══════════════════════════════════════════════════════════════
\section{Casos de Teste}
% ══════════════════════════════════════════════════════════════

\subsection*{Convenção de identificação}

\begin{itemize}
  \item \textbf{CT-XX.Y} $\rightarrow$ Caso de teste número XX, sub-caso Y
  \item Vinculado ao requisito funcional correspondente em \texttt{03-especificacao-requisitos-SRS.md}
\end{itemize}

\subsection{RF-05 --- Cálculo de NDVI}

\renewcommand{\arraystretch}{1.4}
\begin{table}[H]
\centering
\footnotesize
\rowcolors{2}{tabrow}{white}
\begin{tabularx}{\textwidth}{@{}p{1.5cm}p{3.5cm}p{3.5cm}X@{}}
\toprule
\rowcolor{tabheader}\textcolor{white}{\textbf{ID}} & \textcolor{white}{\textbf{Caso}} & \textcolor{white}{\textbf{Entrada}} & \textcolor{white}{\textbf{Resultado Esperado}}\\
\midrule
CT-01.1 & NDVI no intervalo válido & Bandas NIR e Red sintéticas (uint16, 100$\times$100) & Array com $\min \geq -1$ e $\max \leq 1$\\
CT-01.2 & Vegetação saudável produz NDVI alto & NIR alto (3000--8000), Red baixo (500--2000) & $\mathrm{mean} > 0{,}4$\\
CT-01.3 & Fórmula matemática correta & NIR = 8000, Red = 2000 & NDVI = 0{,}6 ($= 6000/10000$) com tolerância $10^{-6}$\\
CT-01.4 & Tratamento de divisão por zero & NIR = 0, Red = 0 & Todos os pixels NaN (não crash)\\
CT-01.5 & NDVI zero quando NIR == Red & NIR = Red = 3000 & NDVI $\approx 0{,}0$\\
\bottomrule
\end{tabularx}
\end{table}

\subsection{RF-06 --- Cálculo de EVI}

\begin{table}[H]
\centering
\footnotesize
\rowcolors{2}{tabrow}{white}
\begin{tabularx}{\textwidth}{@{}p{1.5cm}p{4cm}p{3.5cm}X@{}}
\toprule
\rowcolor{tabheader}\textcolor{white}{\textbf{ID}} & \textcolor{white}{\textbf{Caso}} & \textcolor{white}{\textbf{Entrada}} & \textcolor{white}{\textbf{Resultado Esperado}}\\
\midrule
CT-02.1 & EVI no intervalo válido & NIR, Red, Blue sintéticos & $\min \geq -1$ e $\max \leq 1$ (clipping aplicado)\\
CT-02.2 & Coeficientes padrão & NIR=8000, Red=2000, Blue=1000 & EVI clipado em 1{,}0\\
CT-02.3 & Coeficientes customizados produzem resultado diferente & Mesmas bandas, $G=2{,}0$ vs $G=2{,}5$ & Resultados numericamente diferentes\\
\bottomrule
\end{tabularx}
\end{table}

\subsection{RF-10 --- Estatísticas Descritivas}

\begin{table}[H]
\centering
\footnotesize
\rowcolors{2}{tabrow}{white}
\begin{tabularx}{\textwidth}{@{}p{1.5cm}p{4cm}p{3.5cm}X@{}}
\toprule
\rowcolor{tabheader}\textcolor{white}{\textbf{ID}} & \textcolor{white}{\textbf{Caso}} & \textcolor{white}{\textbf{Entrada}} & \textcolor{white}{\textbf{Resultado Esperado}}\\
\midrule
CT-03.1 & Todas as chaves presentes & Array NDVI 50$\times$50 & Dict contém \texttt{min, max, mean, median, std, count}\\
CT-03.2 & Valores matematicamente corretos & \texttt{[[0.1,0.2,0.3],[0.4,0.5,0.6]]} & min=0{,}1; max=0{,}6; mean=0{,}35; count=6\\
CT-03.3 & NaN ignorados no cálculo & Array 20$\times$20 com 100 NaN & count = 300 ($= 400 - 100$)\\
CT-03.4 & Array totalmente NaN & Array 10$\times$10 todo NaN & count = 0; mean = 0{,}0 (sem crash)\\
\bottomrule
\end{tabularx}
\end{table}

\subsection{RF-11 --- Histograma}

\begin{table}[H]
\centering
\footnotesize
\rowcolors{2}{tabrow}{white}
\begin{tabularx}{\textwidth}{@{}p{1.5cm}p{4cm}p{3.5cm}X@{}}
\toprule
\rowcolor{tabheader}\textcolor{white}{\textbf{ID}} & \textcolor{white}{\textbf{Caso}} & \textcolor{white}{\textbf{Entrada}} & \textcolor{white}{\textbf{Resultado Esperado}}\\
\midrule
CT-04.1 & Padrão 50 bins & Array NDVI saudável & \texttt{len(counts) == 50} e \texttt{len(bins) == 50}\\
CT-04.2 & Bins customizados & \texttt{bins=10} & \texttt{len(counts) == 10}\\
CT-04.3 & Soma de counts = total de pixels & Array 50$\times$50 & \texttt{sum(counts) == 2500}\\
CT-04.4 & Array vazio & Array 5$\times$5 todo NaN & Listas vazias retornadas (sem crash)\\
\bottomrule
\end{tabularx}
\end{table}

\subsection{Casos de Borda (CT-05)}

\begin{table}[H]
\centering
\footnotesize
\rowcolors{2}{tabrow}{white}
\begin{tabularx}{\textwidth}{@{}p{1.5cm}p{4cm}p{3.5cm}X@{}}
\toprule
\rowcolor{tabheader}\textcolor{white}{\textbf{ID}} & \textcolor{white}{\textbf{Caso}} & \textcolor{white}{\textbf{Entrada}} & \textcolor{white}{\textbf{Resultado Esperado}}\\
\midrule
CT-05.1 & Conversão uint16 $\rightarrow$ float32 & NIR e Red como uint16 & \texttt{ndvi.dtype == np.float32}\\
CT-05.2 & NDVI negativo para água & NIR=500, Red=2000 (água) & Todos os pixels $< 0$\\
\bottomrule
\end{tabularx}
\end{table}

\subsection{RF-20 --- Endpoint \texttt{/health}}

\begin{table}[H]
\centering
\footnotesize
\rowcolors{2}{tabrow}{white}
\begin{tabularx}{\textwidth}{@{}p{1.5cm}p{4cm}p{3.5cm}X@{}}
\toprule
\rowcolor{tabheader}\textcolor{white}{\textbf{ID}} & \textcolor{white}{\textbf{Caso}} & \textcolor{white}{\textbf{Requisição}} & \textcolor{white}{\textbf{Resultado Esperado}}\\
\midrule
CT-06.1 & Health responde 200 & \texttt{GET /health} & HTTP 200\\
CT-06.2 & Estrutura da resposta & \texttt{GET /health} & JSON com \texttt{status: "healthy"} e \texttt{products\_dir: bool}\\
\bottomrule
\end{tabularx}
\end{table}

\subsection{Endpoint Raiz, Produtos e Validação}

\begin{table}[H]
\centering
\footnotesize
\rowcolors{2}{tabrow}{white}
\begin{tabularx}{\textwidth}{@{}p{1.5cm}p{4.5cm}p{3.5cm}X@{}}
\toprule
\rowcolor{tabheader}\textcolor{white}{\textbf{ID}} & \textcolor{white}{\textbf{Caso}} & \textcolor{white}{\textbf{Requisição}} & \textcolor{white}{\textbf{Resultado Esperado}}\\
\midrule
CT-07.1 & Raiz retorna info da API & \texttt{GET /} & JSON com \texttt{message}, \texttt{version}, \texttt{endpoints}\\
CT-08.1 & Endpoint \texttt{/products} responde 200 & \texttt{GET /products} & HTTP 200\\
CT-08.2 & Estrutura JSON correta & \texttt{GET /products} & \texttt{products} (lista) e \texttt{count} (int)\\
CT-09.1 & Campos obrigatórios faltando & \texttt{POST /calculate-index} body vazio & HTTP 422\\
CT-09.2 & Tipo de coordenada inválido & longitude=``invalid'' (string) & HTTP 422\\
CT-09.3 & \texttt{index\_type} desconhecido & index\_type=``INVALID\_INDEX'' & HTTP 400/404/500\\
\bottomrule
\end{tabularx}
\end{table}

\subsection{Conversão de Coordenadas (CT-10)}

\begin{table}[H]
\centering
\footnotesize
\rowcolors{2}{tabrow}{white}
\begin{tabularx}{\textwidth}{@{}p{1.5cm}p{4.5cm}p{3.5cm}X@{}}
\toprule
\rowcolor{tabheader}\textcolor{white}{\textbf{ID}} & \textcolor{white}{\textbf{Caso}} & \textcolor{white}{\textbf{Entrada}} & \textcolor{white}{\textbf{Resultado Esperado}}\\
\midrule
CT-10.1 & Polígono Shapely válido & 4 coordenadas formando triângulo fechado & \texttt{polygon.is\_valid == True}, \texttt{area > 0}\\
CT-10.2 & Centroide dentro do bbox & Coordenadas em Picos-PI & $-42{,}6194 \leq c_x \leq -42{,}6180$ e $-4{,}8809 \leq c_y \leq -4{,}8798$\\
\bottomrule
\end{tabularx}
\end{table}

% ══════════════════════════════════════════════════════════════
\section{Matriz de Rastreabilidade Requisito $\times$ Caso de Teste}
% ══════════════════════════════════════════════════════════════

\begin{table}[H]
\centering
\small
\rowcolors{2}{tabrow}{white}
\begin{tabularx}{\textwidth}{@{}lX@{}}
\toprule
\rowcolor{tabheader}\textcolor{white}{\textbf{Requisito (SRS)}} & \textcolor{white}{\textbf{Casos de Teste}}\\
\midrule
RF-04 (Listar produtos) & CT-08.1, CT-08.2\\
RF-05 (Calcular NDVI) & CT-01.1 a CT-01.5\\
RF-06 (Calcular EVI) & CT-02.1, CT-02.2, CT-02.3\\
RF-10 (Estatísticas) & CT-03.1 a CT-03.4\\
RF-11 (Histograma) & CT-04.1 a CT-04.4\\
RF-20 (Health check) & CT-06.1, CT-06.2\\
RNF-04 (Confiabilidade --- erros claros) & CT-01.4, CT-03.4, CT-04.4, CT-09.1, CT-09.2, CT-09.3\\
RNF-05 (Manutenibilidade --- separação de camadas) & Toda a suite (testa \texttt{SentinelProcessor} isoladamente das rotas HTTP)\\
\bottomrule
\end{tabularx}
\end{table}

% ══════════════════════════════════════════════════════════════
\section{Riscos e Mitigações}
% ══════════════════════════════════════════════════════════════

\begin{table}[H]
\centering
\small
\rowcolors{2}{tabrow}{white}
\begin{tabularx}{\textwidth}{@{}XX@{}}
\toprule
\rowcolor{tabheader}\textcolor{white}{\textbf{Risco}} & \textcolor{white}{\textbf{Mitigação}}\\
\midrule
Testes dependem de produtos Sentinel-2 reais ($\sim$1\,GB) & Uso de fixtures sintéticas via NumPy. Apenas testes manuais usam dados reais (na demo).\\
Ambiente de CI sem GDAL/rasterio & Execução em máquina local com GDAL instalado via Homebrew. Não há CI configurado neste MVP (escopo P2).\\
Cobertura limitada de funções de I/O (\texttt{read\_band}, \texttt{crop\_to\_geometry}) & Aceitável para o MVP --- essas funções são thin wrappers sobre rasterio, validados em demonstração manual.\\
\bottomrule
\end{tabularx}
\end{table}

% ══════════════════════════════════════════════════════════════
\section{Critério de Saída}
% ══════════════════════════════════════════════════════════════

\begin{successbox}
A entrega da P2 está pronta quando:
\begin{itemize}
  \item[\checkmark] Todos os 28 casos de teste foram executados
  \item[\checkmark] Taxa de sucesso $\geq 95\%$ \textbf{(atingido: 100\%)}
  \item[\checkmark] Relatório consolidado com evidências está disponível em \texttt{09-relatorio-de-testes.pdf}
  \item[\checkmark] Suite executa em $< 30$\,s \textbf{(atingido: $\sim 17$\,s na primeira execução; $< 1$\,s em \textit{re-runs})}
\end{itemize}
\end{successbox}

\end{document}
